\documentclass{ltjsbook}
\usepackage{docmute} 
\usepackage{../../myStyle} 

\begin{document}
\chapter{行列操作のRコード}
\label{LA_R}
\section{行列の基本変形}
行列の計算は，一度は手計算をして計算方法を「身体化」した方が良いと思いますが，単純な計算でも繰り返しが多くて人間がやるとどうしてもミスが生じます。
そこでRなど計算機を使っての検算をすることになると思います。本書執筆時に，行列の基本変形(セクション\ref{ElementaryTransformation},Pp.\pageref{ElementaryTransformation})に言及しましたが，その操作をするコードを以下に示します。

\begin{lstlisting}[language=C++,label={Pic},caption={ある行に0ではない数字を掛ける}]
pic <- function(size,target,constant){
    vec = rep(1,size)
    vec[target] = constant
    return(diag(vec))
}
\end{lstlisting}

\begin{lstlisting}[language=C++,label={Pijc},caption={ある行に他の行を何倍かしたものを加える}]
pij <- function(size,from,to){
    vec = rep(1,size)
    vec[from] = 0
    vec[to] = 0
    mat = diag(vec)
    mat[from,to] = 1
    return(mat)
}
\end{lstlisting}

\begin{lstlisting}[language=C++,label={Pij},caption={2つの行を入れ替える}]
pij <- function(size,from,to){
    vec = rep(1,size)
    vec[from] = 0
    vec[to] = 0
    mat = diag(vec)
    mat[from,to] = 1
    return(mat)
}
\end{lstlisting}

これらのコードはいずれも，引数として行列のサイズを入れます。関数の中ではRの\verb|diag|関数で単位行列を作ってから，必要な要素だけ入れ替えて返すようにしています。
例題にあった，拡大係数行列
\[ \bm{A} = \begin{pmatrix} 1&2&4&3\\2&1&-1&3\\3&1&4&1\end{pmatrix} \]
を変形する操作ですが，最終的には
\[ \bm{P}_{32}(-2)\bm{P}_{31}(1)\bm{P}_{3}(1/2)\bm{P}_{23}(-7)\bm{P}_{21}(2)\bm{P}_{2}(1/14)\bm{P}_{13}(-3)\bm{P}_{12}(-5)\bm{A} = \bm{I} \]
という計算をすることになったのでした。これを今回のコードを使ってやると，以下のようになります。
\begin{lstlisting}[language=C++,label={P},caption={拡大係数行列の変形}]
A <- matrix(c(1,5,3,-2,4,1,-5,3,-3,3,1,6),ncol=4)
size <- 3
P1 <- pijc(size,2,1,-5)
P2 <- pijc(size,3,1,-3)
P3 <- pic(size,2,1/14)
P4 <- pijc(size,1,2,2)
P5 <- pijc(size,3,2,-7)
P6 <- pic(size,3,-1/2)
P7 <- pijc(size,1,3,1)
P8 <- pijc(size,2,3,-2)
P8 %*% P7 %*% P6 %*% P5 %*% P4 %*% P3 %*% P2 %*% P1 %*% A
\end{lstlisting}

Rの行列の積は\verb|%*%|という記号になることに注意してください。

また，全ての操作を結合した行列$\bm{P}$が逆行列に一致することを確認するために，Rの逆行列関数\verb|solve|を使って算出したものと見比べてみてください(Rの出力\ref{RS:LAout01})。
\begin{lstlisting}[language=C++,label={solve},caption={逆行列の算出}]
I <- diag(rep(1,size))
P8 %*% P7 %*% P6 %*% P5 %*% P4 %*% P3 %*% P2 %*% P1 %*% I
solve(matrix(c(1,5,3,-2,4,1,-5,3,-3),ncol=3))
\end{lstlisting}

\begin{Rscreen}{逆行列の検算結果}{LAout01}
\begin{verbatim}
> P8 %*% P7 %*% P6 %*% P5 %*% P4 %*% P3 %*% P2 %*% P1 %*% I
           [,1]       [,2] [,3]
[1,]  0.5357143  0.3928571 -0.5
[2,] -0.8571429 -0.4285714  1.0
[3,]  0.2500000  0.2500000 -0.5
> solve(matrix(c(1,5,3,-2,4,1,-5,3,-3),ncol=3))
           [,1]       [,2] [,3]
[1,]  0.5357143  0.3928571 -0.5
[2,] -0.8571429 -0.4285714  1.0
[3,]  0.2500000  0.2500000 -0.5
 \end{verbatim}
 \end{Rscreen}

\section{固有値の近似演算}

固有値と固有ベクトルの計算は基本的に機械に任せていればよく，Rでは\verb|eigen|という関数で固有値と固有ベクトルを返してくれます。ここではちょっとマニアックに，どうやって計算しているのかについて考えてみたいと思います。

固有値や固有ベクトルの計算は固有方程式を解くことにあり，変数が多くなると近似的な計算をすることになり，計算機科学の分野ではさまざまな近似計算法が提案されています。ここではそのうちの2つをご紹介します。
\subsection{パワー法}
まずはパワー法から。ここでのパワーpowerとは力という意味ではなく，累乗の方の意味です。

ある行列$\bm{A}$の固有値を$\lambda_1,\lambda_2,...,\lambda_n$とします。かつ，これらが絶対値の大きい順に並んでいるとしましょう。すなわち，
\[ \lambda_1 > \lambda_2 > \cdots > \lambda_n \]
です。またこの各固有値に対応する固有ベクトルをそれぞれ$\bm{v}_1,\bm{v}_2,\cdots,\bm{v}_n$と表します。任意のベクトル$\bm{v}$は$\bm{v}_1,\bm{v}_2,\cdots,\bm{v}_n$の線形結合で表すことができるという性質を利用して，
\[ \bm{v} = c_1\bm{v}_1 + c_2\bm{v}_2 + \cdots + c_n\bm{v}_n \]
と書くことができます。このとき，$c_j$はあるスカラーです。さて，この$\bm{v}$を行列$\bm{A}$にかけるとどうなるかというと，
\[
    \begin{array}{lll}
        \bm{Av} &=& \bm{A}(c_1\bm{v}_1) + \bm{A}(c_2\bm{v}_2) + \cdots + \bm{A}(c_n\bm{v}_n)\\
        \\
        &=& c_1\bm{Av}_1 + c_2\bm{Av}_2 + \cdots + c_n\bm{Av}_n\\
        \\
        &&\mbox{$\bm{v}$はいずれも固有ベクトルだから}\\
        \\
        &=& c_1\lambda_1\bm{v}_1 + c_2 \lambda_2\bm{v}_2 + \cdots + c_n\lambda_n\bm{v}_n
    \end{array}
\]
ここで$\lambda_1$を括り出すと，
\[ = \lambda_1(c_1\bm{v}_1 + c_2(\lambda_2/\lambda_1)\bm{v}_2 + \cdots + c_n(\lambda_n/\lambda_1)\bm{v}_n ) \]
となります。このようにしながら，$\bm{A}$を累乗していく，すなわち$\bm{A}\bm{A}\bm{A}\cdots\bm{Av}$としていくと，
\[ \bm{A}^p\bm{v} = \lambda_1^pc_1\bm{v}_1 + c_2(\lambda_2/\lambda_1)^p\bm{v}_2 + \cdots + c_n(\lambda_n/\lambda_1)^p\bm{v}_n \]
となっていきます。ここで$p$回る以上を繰り返しているわけですが，$\lambda1 > \lambda_2$のように最初の固有値$\lambda_1$は最大のものですから，$c_2(\lambda_2/\lambda_1)^p \bm{v}_2$は$(\lambda_2/\lambda_1)^p$がどんどん0に近づいていきます。以下の$c_n(\lambda_n/\lambda_1)^p \bm{v}_n$についても同様で，繰り返し続けると
\[ \bm{A}^p\bm{v} \approx \lambda_1^p c_1 \bm{v}_1 \]
となることがわかります。これを利用して，固有値を求める式は
\[ \lambda_1 = \lim_{p \to \infty}\frac{(\bm{A}^{p+1}\bm{v})_r}{(\bm{A}^p\bm{v})_r} \]

とすることができます。式中の$r$は分子・分母のベクトルの，第$r$番目の要素を指しています。

これがパワー法による固有値の求め方です。この方法では最大固有値しか求められませんが，第2固有値は第1固有値の要素を元の行列から抜き去った行列，すなわち$\bm{A} - \bm{v}_1 \bm{v}_1'$を作り，第1固有値の次元情報が抜けた行列の中の最大固有値，というようにして求めていくことができます。

数値例で見てみましょう。

$4\times4$の行列$\bm{A}$を次のような行列とします。
\[ \bm{A}= \begin{pmatrix} 1.0 & 0.3 & 0.4 & 0.9\\ 0.3 & 1.0 & 0.5 & 0.4\\  0.4 & 0.5 & 1.0 & 0.7\\ 0.9 & 0.4 & 0.7 & 1.0 \end{pmatrix} \]
これに，まずあるベクトル$\bm{v} = \begin{pmatrix} 1.0\\ 0.0\\ 0.0\\ 0.0 \end{pmatrix}$をかけると,

\[ \bm{A}\bm{v}=\begin{pmatrix} 1.0 & 0.3 & 0.4 & 0.9\\ 0.3 & 1.0 & 0.5 & 0.4\\  0.4 & 0.5 & 1.0 & 0.7\\ 0.9 & 0.4 & 0.7 & 1.0 \end{pmatrix} \begin{pmatrix} 1.0\\ 0.0\\ 0.0\\ 0.0 \end{pmatrix}= \begin{pmatrix} 1.0\\ 0.3\\ 0.4\\ 0.9 \end{pmatrix}\]
この新しくできたベクトルを$\bm{v}\dagger$として，もう一度$\bm{A}$にかけると
\[ \bm{A}\bm{v}\dagger=\begin{pmatrix} 1.0 & 0.3 & 0.4 & 0.9\\ 0.3 & 1.0 & 0.5 & 0.4\\  0.4 & 0.5 & 1.0 & 0.7\\ 0.9 & 0.4 & 0.7 & 1.0 \end{pmatrix} \begin{pmatrix} 1.0\\ 0.3\\ 0.4\\ 0.9 \end{pmatrix}=\begin{pmatrix} 2.06\\ 1.16\\ 1.58\\ 2.20 \end{pmatrix} \]
となります。この新しくできたベクトル$\bm{v}\dagger\dagger$の要素は大きいものでは2.0を超えているので，最大値$2.2$を取り出して少しスケールを整えてみましょう。
\[ \bm{w} = 2.2 \begin{pmatrix} 0.936\\ 0.527\\ 0.718\\ 1.000 \end{pmatrix}\]
この大きさの整えられたベクトル$\bm{w}$を次の計算に用いることにします。
\[ \bm{A}\bm{w}=\begin{pmatrix} 1.0 & 0.3 & 0.4 & 0.9\\ 0.3 & 1.0 & 0.5 & 0.4\\  0.4 & 0.5 & 1.0 & 0.7\\ 0.9 & 0.4 & 0.7 & 1.0 \end{pmatrix} \begin{pmatrix} 0.936\\ 0.527\\ 0.718\\ 1.00 \end{pmatrix}=
\begin{pmatrix} 2.282\\ 1.567\\ 2.056\\ 2.556 \end{pmatrix} \]
ここでも同様に，新しくできたベクトルは最大要素で割ってサイズを整え，$\bm{w}\dagger = 2.556 \begin{pmatrix} 0.893\\ 0.613\\ 0.804\\ 1.000 \end{pmatrix} $として次の計算に用いることにします。

以下のプロセスは同様なので，数値の変化だけ追いかけていきましょう。
$$\bm{A}\bm{w}\dagger=\left( \begin{array}{cccc} 1.0 & 0.3 & 0.4 & 0.9\\ 0.3 & 1.0 & 0.5 & 0.4\\  0.4 & 0.5 & 1.0 & 0.7\\ 0.9 & 0.4 & 0.7 & 1.0 \end{array} \right) \left( \begin{array}{c} 0.893\\ 0.613\\ 0.804\\ 1.00 \end{array} \right)=\left( \begin{array}{c} 2.298\\ 1.683\\ 2.168\\ 2.612 \end{array} \right)=2.612\left( \begin{array}{c} 0.880\\ 0.644\\ 0.830\\ 1.00 \end{array} \right)$$
$$\to\left( \begin{array}{cccc} 1.0 & 0.3 & 0.4 & 0.9\\ 0.3 & 1.0 & 0.5 & 0.4\\  0.4 & 0.5 & 1.0 & 0.7\\ 0.9 & 0.4 & 0.7 & 1.0 \end{array} \right) \left( \begin{array}{c} 0.880\\ 0.644\\ 0.830\\ 1.00 \end{array} \right)=\left( \begin{array}{c} 2.305\\ 1.724\\ 2.204\\ 2.631 \end{array} \right)=2.631\left( \begin{array}{c} 0.876\\ 0.655\\ 0.838\\ 1.00 \end{array} \right)$$
$$\to\left( \begin{array}{cccc} 1.0 & 0.3 & 0.4 & 0.9\\ 0.3 & 1.0 & 0.5 & 0.4\\  0.4 & 0.5 & 1.0 & 0.7\\ 0.9 & 0.4 & 0.7 & 1.0 \end{array} \right) \left( \begin{array}{c} 0.876\\ 0.655\\ 0.838\\ 1.00 \end{array} \right)=\left( \begin{array}{c} 2.308\\ 1.737\\ 2.216\\ 2.637 \end{array} \right)=2.637\left( \begin{array}{c} 0.875\\ 0.659\\ 0.840\\ 1.00 \end{array} \right)$$
$$\to\left( \begin{array}{cccc} 1.0 & 0.3 & 0.4 & 0.9\\ 0.3 & 1.0 & 0.5 & 0.4\\  0.4 & 0.5 & 1.0 & 0.7\\ 0.9 & 0.4 & 0.7 & 1.0 \end{array} \right) \left( \begin{array}{c} 0.875\\ 0.659\\ 0.840\\ 1.00 \end{array} \right)=\left( \begin{array}{c} 2.309\\ 1.741\\ 2.220\\ 2.639 \end{array} \right)=2.639\left( \begin{array}{c} 0.875\\ 0.660\\ 0.841\\ 1.00 \end{array} \right)$$
$$\to\left( \begin{array}{cccc} 1.0 & 0.3 & 0.4 & 0.9\\ 0.3 & 1.0 & 0.5 & 0.4\\  0.4 & 0.5 & 1.0 & 0.7\\ 0.9 & 0.4 & 0.7 & 1.0 \end{array} \right) \left( \begin{array}{c} 0.875\\ 0.660\\ 0.841\\ 1.00 \end{array} \right)=\left( \begin{array}{c} 2.309\\ 1.743\\ 2.221\\ 2.640 \end{array} \right)=2.640\left( \begin{array}{c} 0.875\\ 0.660\\ 0.841\\ 1.00 \end{array} \right)$$

このように，取り出される数値が徐々に一定の値(今回は$\to2.640$)に収束していくし，ベクトルの方も一定の値に収まっていくことになります。この変化量がとても小さくなると「十分に収束した」とされ，取り出された値＝固有値，ベクトル＝固有ベクトル，とすることができます。

この計算をするRコードを添付します。
\begin{lstlisting}[language=C++,label={Power},caption={パワー法による最大固有値の算出関数}]
power <- function(A, max.itr = 100, eps = 1e-20) {
  u <- rep(1, length = nrow(A))
  v <- rep(0, length = nrow(A))
  e.val <- 0
  for (i in 1:max.itr) {
    v <- A %*% u
    if (abs(max(v) - e.val) < eps) {
      break
    }
    e.val <- max(v)
    e.vec <- u<- v / e.val
  }
  return(list(Eval = e.val, Evec=e.vec))
}
\end{lstlisting}

\subsection{QR法}
行列を$\bm{A}=\bm{QR}$の形に分解することをQR分解といいます。ここで$\bm{Q}$は直交行列，$\bm{R}$は上三角行列です。上三角行列というのは，非ゼロ要素が体格の右上半分だけにある行列，すなわち
\[
    \begin{pmatrix}
        * & * & * & \cdots & * & * \\
        0 & * & * & \cdots & * & * \\
        0 & 0 & * & \cdots & * & * \\
        0 & 0 & 0 & \cdots & * & * \\
          &   &  \cdots & & \\
        0 & 0 & 0 & \dots & * & * \\
        0 & 0 & 0 & \dots & 0 & * \\
    \end{pmatrix}
\]
という形の行列です。ここで*には任意の数値です。

この分解を使って，
\begin{enumerate}
	\item $\bm{A}$を直交行列$\bm{Q}$と上三角行列$\bm{R}$の積に分解する。すなわち，$\bm{A}^{(t)} = \bm{Q} \bm{R}$
	\item その結果を逆順に掛ける。$\bm{A}^{(t+1)}=\bm{R} \bm{Q}$
\end{enumerate}
という2つの操作を反復して，固有値を求めることになります。これらの計算は，固有値を変化させずに行列の形を変形するアルゴリズムであり，このような反復を繰り返すと非対角要素の絶対値は次第に小さくなり，$t$が十分に大きくなれば$\bm{A}^{(t)}$は対角行列に収束し，その対角項が固有値となるという技です。

$\bm{A}$を$\bm{Q}$と$\bm{R}$との積に分解する方法はいくつかありますが，ひとつはセクション\ref{GMorth}，Pp.\pageref{GMorth}で解説したグラム・シュミットの正規直交化法が使えます。

この方法では，行列$\bm{A}$の第一列$\bm{a}_1$，第二列$\bm{a}_2$，$\cdots$，第$n$列$\bm{a}_n$を元に，行列$\bm{Q}$の第一列$\bm{q}_1$，第二列$\bm{q}_2$，$\cdots$，第$n$列$\bm{q}_n$を作っていきます。このとき$\bm{Q}$は
\begin{itemize}
	\item $\bm{q}_i$と$\bm{q}_j$が互いに直交する。すなわち，$(\bm{q}_i,\bm{q}_j)=0.$ただし$i \neq j$
	\item 各$\bm{q}_i$の長さは$1$である。すなわち，$(\bm{q}_i,\bm{q}_i)=1$
\end{itemize}
の条件を満たすように作っていきます。

そこで，
\begin{equation}
\begin{array}{ccl}
\bm{b}_1 & = & \bm{a}_1\\
\bm{b}_2 & = & (\bm{a}_2 から\bm{a}_1 に平行な成分を抜き取ったもの)\\
\bm{b}_3 & = & (\bm{a}_3 から\bm{a}_1,\bm{a}_2に平行な成分を抜き取ったもの)\\
& \cdots & \\
\bm{b}_n & = & (\bm{a}_n から\bm{a}_1,\cdots \bm{a}_{n-1}に平行な成分を抜き取ったもの)\\
\end{array}
\end{equation}
という要領で，$\bm{b}_1,\bm{b}_2 \,\cdots \bm{b}_n$を作り，各$\bm{b}_i$をその長さ$\| \bm{b}_i \|$で割って$\bm{q}_i$にすることになります。

数値例を見てみましょう。行列$\bm{A}$を$\bm{QR}$の形に分解してみます。
\[
    \begin{pmatrix}
        1.0 & 0.3 & 0.4 & 0.9 \\ 
        0.3 & 1.0 & 0.5 & 0.4 \\ 
        0.4 & 0.5 & 1.0 & 0.7 \\ 
        0.9 & 0.4 & 0.7 & 1.0 \end{pmatrix}
        =\\
    \begin{pmatrix}
        0.697 & -0.286 & -0.338 & -0.565 \\ 
        0.209 & 0.903 & -0.374 & 0.025 \\ 
        0.279 & 0.299 & 0.855 & -0.319 \\ 
        0.627 & -0.116 & 0.120 & 0.761 \\ 
    \end{pmatrix}
    \begin{pmatrix}
        1.435 & 0.808 & 1.101 & 1.533 \\ 
        0.000 & 0.920 & 0.555 & 0.197 \\ 
        0.000 & 0.000 & 0.617 & 0.265 \\ 
        0.000 & 0.000 & 0.000 & 0.039 \\ 
    \end{pmatrix}
\]
そしてこれから，$\bm{A}^{(1)} = \bm{RQ}$を作ります。
\[
    \begin{pmatrix}
        2.437 & 0.470 & 0.338 & 0.025 \\ 
        0.470 & 0.974 & 0.153 & -0.005 \\ 
        0.338 & 0.153 & 0.559 & 0.005 \\ 
        0.025 & -0.005 & 0.005 & 0.030 \\ 
    \end{pmatrix}
\]
この操作を繰り返して，$\bm{A}=\bm{QR} \to \bm{RQ} = \bm{A}^{(1)} \to \bm{A}^{(1)} = \bm{QR} \to \cdots$を続けてみたいと思います。
\[
    \bm{A}^{(2)} = 
    \begin{pmatrix}
        2.620 & 0.180 & 0.067 & 0.00\\
        0.180 & 0.859 & 0.037 & 0.00\\
        0.067 & 0.037 & 0.492 & 0.00\\
        0.000 & 0.000 & 0.000 & 0.03\\
    \end{pmatrix}
\]

\[
    \bm{A}^{(3)} = 
    \begin{pmatrix}
        2.638 & 0.059 & 0.012 & 0.00\\
        0.059 & 0.844 & 0.018 & 0.00\\
        0.012 & 0.018 & 0.488 & 0.00\\
        0.000 & 0.000 & 0.000 & 0.03\\
    \end{pmatrix}
\]

\[
    \bm{A}^{(4)} = 
    \begin{pmatrix}
        2.640 & 0.019 & 0.002 & 0.00\\
        0.019 & 0.843 & 0.010 & 0.00\\
        0.002 & 0.010 & 0.487 & 0.00\\
        0.000 & 0.000 & 0.000 & 0.03\\
    \end{pmatrix}
\]

\[
    \bm{A}^{(5)} = 
    \begin{pmatrix}
        2.640 & 0.006 & 0.000 & 0.00\\
        0.006 & 0.843 & 0.006 & 0.00\\
        0.000 & 0.006 & 0.487 & 0.00\\
        0.000 & 0.000 & 0.000 & 0.03\\
    \end{pmatrix}
\]
どんどん対角化されていくのが見えてきましたね。これを10回もすると，
\[
    \bm{A}^{(10)} = 
    \begin{pmatrix}
        2.64 & 0.000 & 0.000 & 0.00\\
        0.00 & 0.843 & 0.000 & 0.00\\
        0.00 & 0.000 & 0.487 & 0.00\\
        0.00 & 0.000 & 0.000 & 0.03\\
            \end{pmatrix}
\]
となり，対角項に4つの固有値が並んでいるのが見えると思います。


最後に正方行列$\bm{A}$を与えると，QR分解するRコードを添付します。
\begin{lstlisting}[language=C++,label={QRdecompose},caption={QR分解するRコード}]
QRdecompose <- function(A){
    Q = matrix(0,ncol=ncol(A),nrow=nrow(A))
    R = matrix(0,ncol=ncol(A),nrow=nrow(A))
    for(i in 1:N){
        b <- A[,i]
        for(j in 1:i){
            if(j<i){
                R[j,i] <- t(A[,i])%*%Q[,j]
                b <- b -  R[j,i]*Q[,j]
            }
        }
        R[i,i] <- sqrt(sum(b^2))
        Q[,i] <- b/R[i,i]
    }
    return(list(Q=Q,R=R))
}
\end{lstlisting}

\end{document}